\begin{ex}
Cho tập hợp $A=\{x\in\mathbb{N}|4<x\le9\}$. Tập hợp A viết dưới dạng liệt kê là:
\choice
{\{4; 5; 6; 7; 8; 9\}}
{\{4; 5; 6; 7; 8\}}
{\{5; 6; 7; 8\}}
{\True \{5; 6; 7; 8; 9\}}
\loigiai{
Đáp án đúng: D
}
\end{ex}

\begin{ex}
Số tự nhiên $x$ thỏa mãn $(x-29)-11=0$ là:
\choice
{30}
{39}
{\True 40}
{41}
\loigiai{
Đáp án đúng: C
}
\end{ex}

\begin{ex}
Kết quả phép tính $312 - 97$ bằng:
\choice
{212}
{\True 215}
{225}
{409}
\loigiai{
Đáp án đúng: B
}
\end{ex}

\begin{ex}
Tích $2525:25.5-5$ là:
\choice
{50}
{0}
{\True 500}
{Kết quả khác}
\loigiai{
Đáp án đúng: C
}
\end{ex}

\begin{ex}
Thứ tự thực hiện các phép tính đối với biểu thức không có dấu ngoặc là:
\choice
{Nhân và chia $\rightarrow$ Lũy thừa $\rightarrow$ Cộng và trừ}
{\True Lũy thừa $\rightarrow$ Nhân và chia $\rightarrow$ Cộng và trừ}
{Cộng và trừ $\rightarrow$ Nhân và chia $\rightarrow$ Lũy thừa}
{Lũy thừa $\rightarrow$ Cộng và trừ $\rightarrow$ Nhân và chia}
\loigiai{
Đáp án đúng: B
}
\end{ex}

\begin{ex}
Tính giá trị của lũy thừa $2^{6}$ ta được kết quả là:
\choice
{16}
{32}
{\True 64}
{12}
\loigiai{
Đáp án đúng: C
}
\end{ex}

\begin{ex}
Kết quả phép tính $5^{9} \cdot 5^{5}$ bằng:
\choice
{$5^{45}$}
{\True $5^{14}$}
{$25^{14}$}
{$10^{14}$}
\loigiai{
Đáp án đúng: B
}
\end{ex}

\begin{ex}
Tìm số tự nhiên $x$, biết: $(x - 25) \cdot 2021=0$ ta được kết quả $x$ bằng:
\choice
{0}
{45}
{\True 25}
{26}
\loigiai{
Đáp án đúng: C
}
\end{ex}

\begin{ex}
Thương $q$ và số dư $r$ trong phép chia $a=713$ cho $b=51$ là:
\choice
{$q=14; r=0$}
{\True $q=13; r=50$}
{$q=50; r=13$}
{$q=14; r=1$}
\loigiai{
Đáp án đúng: B
}
\end{ex}

\begin{ex}
Dạng tổng quát của số tự nhiên chia hết cho 3 là?
\choice
{\True $3k (k\in\mathbb{N})$}
{$5k+3 (k\in\mathbb{N})$}
{$3k+1 (k\in\mathbb{N})$}
{$3k+2 (k\in\mathbb{N})$}
\loigiai{
Đáp án đúng: A
}
\end{ex}

\begin{ex}
Một tàu cần chở 1200 hành khách. Mỗi toa có 12 khoang, mỗi khoang có 8 chỗ ngồi. Hỏi tàu hỏa cần ít nhất bao nhiêu toa để chở hết chỗ hành khách đó?
\choice
{12}
{\True 13}
{11}
{14}
\loigiai{
Đáp án đúng: B
}
\end{ex}

\begin{ex}
Trong các số sau, số nào là ước của 10?
\choice
{\True 10}
{0}
{3}
{9}
\loigiai{
Đáp án đúng: A
}
\end{ex}

\begin{ex}
Tổng tất cả các số tự nhiên $x$ thỏa mãn $15 \vdots (2x+1)$ là:
\choice
{\True 10}
{4}
{7}
{6}
\loigiai{
Đáp án đúng: A
}
\end{ex}

\begin{ex}
Tổng (hoặc hiệu) nào dưới đây chia hết cho 7?
\choice
{\True $147+21$}
{$28+16$}
{$49-24$}
{$35-12$}
\loigiai{
Đáp án đúng: A
}
\end{ex}

\begin{ex}
Giá trị tương ứng trong hệ thập phân của số La Mã XXIV là:
\choice
{101015}
{10104}
{\True 24}
{26}
\loigiai{
Đáp án đúng: C
}
\end{ex}

\begin{ex}
Tìm tập hợp các chữ số điền vào dấu * để được số $\overline{14*}$ chia hết cho 2:
\choice
{\True $\{0;2;4;6;8\}$}
{$\{2;4;6;8\}$}
{$\{2\}$}
{$\{0;2;4;6;8;10...\}$}
\loigiai{
Đáp án đúng: A
}
\end{ex}

\begin{ex}
Tất cả các số nguyên tố có một chữ số là:
\choice
{1; 3; 5; 7}
{3; 5; 7}
{2; 3; 5; 9}
{\True 2; 3; 5; 7}
\loigiai{
Đáp án đúng: D
}
\end{ex}

\begin{ex}
Có bao nhiêu số tự nhiên $x$ để $\overline{3x}$ là hợp số?
\choice
{2}
{\True 8}
{7}
{9}
\loigiai{
Đáp án đúng: B
}
\end{ex}

\begin{ex}
Phân tích số 18 ra thừa số nguyên tố. Cách viết nào sau đây là đúng?
\choice
{$18=18.1$}
{$18=10+8$}
{\True $18=2.3^{2}$}
{$18=6+6+6$}
\loigiai{
Đáp án đúng: C
}
\end{ex}

\begin{ex}
Số $126=2.3^{2}.7$ có bao nhiêu ước nguyên tố?
\choice
{\True 3}
{4}
{10}
{12}
\loigiai{
Đáp án đúng: A
}
\end{ex}
